\chapter{补充推导与设计辨析}
\label{ch:deriv}

本章给出第\ref{ch:theory}章主要命题与定理的证明要点，并补充实现层面的设计辨析。表述以严谨性为目标，同时保持与主线符号一致：$|\mathcal{A}|$ 为方形 QAM 星座基数，$X_1$ 为期望用户符号。

\section{命题\ref{prop:mld-equiv}：MLD 与对数损失等价}

\subsection{严格真评分}
对离散标签 $X_1\in\mathcal{A}$，设模型软分布
\begin{equation}
q_a(y)=\frac{\exp(g_a(y))}{\sum_{b\in\mathcal{A}}\exp(g_b(y))}.
\end{equation}
期望交叉熵可分解为
\begin{equation}
\mathbb{E}\big[-\log q_{X_1}(Y)\big]
=\mathbb{E}\big[D_{\mathrm{KL}}\big(f^*(Y)\,\|\,q(Y)\big)\big]
+\mathbb{E}\big[H\big(f^*(Y)\big)\big],
\label{eq:ce-decomp}
\end{equation}
其中 $H(f^*)$ 为后验熵，与模型无关。由 KL 散度非负性，$D_{\mathrm{KL}}(f^*\|q)=0$ 当且仅当 $q=f^*$ a.s.。因此风险 $R(g)$ 的极小化子恰为真后验（logits 可相差与 $a$ 无关的加性函数）。MAP 判决 $\arg\max_a f_a^*=\arg\max_a q_a^*$，即命题\ref{prop:mld-equiv}。

\subsection{与显式 MLD 的联系}
在已知 $H,N_0$ 时，贝叶斯定理给出
\begin{equation}
f_a^*(y)=\frac{p(y\mid X_1=a)}{\sum_b p(y\mid X_1=b)},
\end{equation}
而
\begin{equation}
p(y\mid X_1=a)
=\sum_{x_I\in\mathcal{A}^{K-1}}
p(y\mid X=(a,x_I))\,p(x_I),
\end{equation}
代入复高斯似然即得式\eqref{eq:fstar}。故“穷举边际化”与“风险极小化子”描述同一对象；差别仅在于前者显式用 $H$，后者用样本隐式编码 $H$。

\section{命题\ref{prop:rkhs-approx}：后验的 RKHS 结构}

将联合符号记为 $x\in\mathcal{A}^K$。线性模型下
\begin{equation}
p(y\mid x)
\propto
\exp\Big(
-\frac{\|y-Hx\|^2}{N_0}
\Big)
=
\exp\Big(
-\frac{\|y\|^2}{N_0}
\Big)
\exp\Big(
\frac{2}{N_0}\Re(y^H Hx)
-\frac{\|Hx\|^2}{N_0}
\Big).
\end{equation}
对固定 $a$ 边际化 $x_1=a$ 的 $x$，得到式\eqref{eq:fstar-exp}：分子是有限个形如
$\exp\big(2\Re(y^H \mu_x)/N_0\big)$ 的函数之和（$\mu_x=Hx$）。该类函数落在高斯核特征映射所张成空间的闭包中；再由 Steinwart 通用核定理\cite{steinwart2001} 与 Stone--Weierstrass 型论证，可得一致逼近序列。更细致的范数界依赖带宽与 SNR 的匹配（见下节 $B_{\mathrm{tot}}$）。

\section{定理\ref{thm:representer}：表示定理要点}

标准表示定理\cite{scholkopf2001}：若目标形如
\begin{equation}
\frac1n\sum_i \ell\big(g(x_i),y_i\big)+\lambda\|g\|_{\mathcal{H}}^2,
\end{equation}
且 $\ell$ 只通过有限个点赋值 $\{g(x_i)\}$ 依赖 $g$，则最优解必落在
$\mathrm{span}\{k(\cdot,x_i)\}_{i=1}^n$ 中。对本问题，$g$ 为 $|\mathcal{A}|$ 个分量，正则为 $\sum_a\|g_a\|^2$，损失为 softmax 交叉熵，仅依赖 $\{g_a(y_i)\}$，故逐类展开
$g_a=\sum_i\alpha_i^{(a)}k(\cdot,y_i)$。代入即得有限维问题\eqref{eq:finite-erm}。凸性：交叉熵对 logits 凸，RKHS 正则对 $\alpha$ 在 $K\succeq 0$ 时为凸二次型，故整体对 $\alpha$ 凸（固定核时）。

\section{定理\ref{thm:erm-rate}：excess-risk 与 $n\gtrsim|\mathcal{A}|K/\varepsilon^2$}

\subsection{Rademacher 复杂度路线}
对有界 Lipschitz 损失，标准 excess-risk 界\cite{bartlett2002} 形如
\begin{equation}
R(\hat g)-R(g^*)
\le
2\mathfrak{R}_n(\mathcal{F})
+O\Big(\sqrt{\frac{\ln(1/\delta)}{n}}\Big),
\end{equation}
其中 $\mathfrak{R}_n$ 为假设类的经验 Rademacher 复杂度。对 RKHS 球 $\{g:\|g\|_{\mathcal{H}}\le B\}$ 且 $k(x,x)\le 1$，有
$\mathfrak{R}_n\le B/\sqrt{n}$。对 $|\mathcal{A}|$ 个分量取并，得式\eqref{eq:excess} 中 $B_{\mathrm{tot}}$ 项。

\subsection{$B_{\mathrm{tot}}^2=\mathcal{O}(|\mathcal{A}|K)$ 的量级论证}
当核带宽与噪声同阶时，式\eqref{eq:fstar-exp} 中指数项的“有效支撑”随干扰维度与星座点数目增长。粗略地，每个后验分量的 RKHS 范数随“被边际化的干扰--星座组合”的有效复杂度上升；对 $|\mathcal{A}|$ 个分量求和后得到
$B_{\mathrm{tot}}^2=\mathcal{O}(|\mathcal{A}|K)$ 的工作量级（中期报告中的 $\mathcal{O}(4^Q K)$ 与此同族，此处统一写成 $|\mathcal{A}|$）。令式\eqref{eq:excess} 右端 $\le\varepsilon$，即得
$n\gtrsim B_{\mathrm{tot}}^2/\varepsilon^2\sim|\mathcal{A}|K/\varepsilon^2$。

\subsection{紧致性注记}
各 $f_a^*$ 之和为 1，且由同一 $p(y,x)$ 生成，独立求和 Rademacher 界可能偏松。利用多任务共享结构或 Fano 型信息论下界收紧样本复杂度，仍是后续理论工作。

\section{定理\ref{thm:plaf-gap}：PLAF 与 MMSE 差距的要点}

\subsection{回归目标与 MAP 的错位}
BPSK/QPSK 下判决边界为单一超平面，超平带约束与符号符号一侧裕度一致。高阶 QAM 下实/虚部为多电平，最近邻区域由多条阈值界定；约束 $|g(y)-x|\le\varepsilon$ 不再等价于“落在正确 Voronoi 胞腔”，而是幅度回归。于是学习目标与 0-1/MAP 风险脱节。

\subsection{下界不随 SNR 消失}
对有限样本核回归，标准 excess 项含 $B_g^2/T$。若目标函数的 RKHS 范数随 SNR 与 $|\mathcal{A}|$ 增大（后验更陡、星座更密），则 SER 下界可随 SNR 上升，形成“地板上抬”。而过定系统中 MMSE+切片的 SER 上界随 SNR 指数下降\cite{proakis2008}。两相比较，存在 $\mathrm{SNR}^\star$ 使此后差距为正且扩大——这即定理\ref{thm:plaf-gap} 的逻辑骨架。欠定 $M<K$ 时 MMSE 自身亦有干扰地板，但回归型核方法仍可能叠加额外的样本依赖项。

\section{GA 软分、稳健加载与坐标选择}

在干扰高斯近似下，似然涉及二次型
\begin{equation}
\big(y-a\hat h_1\big)^H R^{-1}\big(y-a\hat h_1\big).
\end{equation}
对不同 $a$ 计算该量（或等价对数分）即得 $z_a\in\mathbb{R}^{|\mathcal{A}|}$。稳健化只替换 $R\to R_{\mathrm{rob}}$，不改变“$|\mathcal{A}|$ 维软分作特征”的结构。由此，$z_{\mathrm{rob}}$ 可视为\textbf{把 CSI 不确定度折叠进充分统计}的一步，使命题\ref{prop:rkhs-approx} 的存在性在有限样本下可数值实现。

\section{Adaptive-MKL 正则 $\Omega$ 的含义}

式\eqref{eq:mkl-reg} 中 $1/\eta_m$ 因子来自对组合核 RKHS 范数的变分表示：在约束 $\sum_m\eta_m=1$、$\eta\ge 0$ 下，
\begin{equation}
\|g\|_{\mathcal{H}_{k_\eta}}^2
=\inf\Big\{
\sum_m\frac{\|g_m\|_{\mathcal{H}_m}^2}{\eta_m}
:\ g=\sum_m g_m
\Big\}.
\end{equation}
因而最小化“损失 $+\lambda\sum_m\|g_m\|^2/\eta_m$”等价于在某个（未知的）组合核 RKHS 中做标准 Tikhonov 正则，同时把核选择交给 $\eta$。熵项 $-\mu H(\eta)$ 则在单纯形内施加温和的均匀先验，避免早期训练把全部质量压到条件数最差的窄核上。

\section{为何 plugin 教师不可过重}

$\hat p=\mathrm{softmax}(z_{\mathrm{rob}})$ 是错误指定模型下的伪后验：既有 GA 近似误差，又有 $\hat H\neq H$ 误差。以 $\hat p$ 为强教师等于要求 $g$ 拟合一个有偏目标，可能增大与真 $X_1$ 相关的 hard 风险。故可部署默认以 hard 为主，plugin 仅作轻正则。

\section{聚合作为同假设类上的模型平均}

设候选 logits $\{L^{(j)}\}$ 均来自同一 $K_\eta$ 的不同 $\alpha^{(j)}$。验证集权重 $w_j\propto \exp(-\beta J_{\mathrm{val}}^{(j)})$，得 $L=\sum_w L^{(j)}$，再解
\begin{equation}
\min_\alpha \|K\alpha^\top-L\|_F^2+\lambda\|\alpha\|^2_K
\end{equation}
投影回表示定理形式。这保证部署期仍只存一份 $(\eta,\alpha)$，且不离开定理\ref{thm:representer} 的假设类。

\section{堆叠的信息论直觉}

一阶段输出 $\mathrm{softmax}(L_1)$ 提供对后验的初步估计；与 $z_{\mathrm{rob}}$ 拼接后，二阶段 RKHS 学习“校正映射”。这类似残差学习，但特征仍显式包含物理统计量。仅在验证 BER 改善且 $\mathrm{SNR}<12\,\mathrm{dB}$ 时启用，以避免高 SNR 稀有错误被二阶段过拟合。

\section{与 MMSE 的关系（再述）}

MMSE 估计是在线性约束下最小化均方误差，再切片。RKHS 方法直接优化分类型 $J$，假设类远大于线性函数。当 $z_{\mathrm{rob}}$ 已含近似白化匹配信息时，核方法是在\textbf{非线性后处理}这些软分，从而在低中 SNR 及模型失配场景获得增益。Oracle 显示：若软分基于真 $H$，非线性后处理可把性能推到 MLD；若基于 $\hat H$，则增益有限但可观。这与定理\ref{thm:plaf-gap} 并不矛盾：本方法优化的是后验分类风险，而非幅度回归。

\section{失败模式清单}

\begin{enumerate}[leftmargin=2em]
  \item 在原始 $Y$ 上拟合尖锐 $f^*$：带宽难选，高 SNR 崩（坐标问题）。
  \item 把高阶 QAM 写成超平带幅度回归：被 MMSE 锁定（定理\ref{thm:plaf-gap}）。
  \item 把 $\hat H$ 代入非稳健 $R$：高 SNR 过自信。
  \item 重 plugin：有偏教师。
  \item 过高 SNR 上过度 L-BFGS / 过多基核：训练数值崩溃。
  \item 用 Oracle 特权信息报告为可部署性能：不可复现于真实链路。
\end{enumerate}
